\section{Conclusion}
\label{sec:conclusion}

We conducted a systematic comparison of induction head behavior on random versus naturalistic data in \gpttwo.
Our results show that induction heads are \emph{not} data-type-specific: the same heads identified by random-sequence detection also show the strongest induction signals on natural text (Spearman $\rho = 0.84$).
However, the nature of their operation differs dramatically across conditions.
Induction attention patterns are 10--40$\times$ weaker on natural text, while late-layer OV copying circuits are up to 29$\times$ \emph{stronger}---revealing a distribution-sensitive form of pattern completion that the standard detection protocol underestimates.

The answer to our motivating question---``are there hidden induction heads that only activate on naturalistic data?''---is: not as separate heads, but as separate \emph{behaviors} of the same heads.
Standard random-sequence detection identifies the right circuits but provides an incomplete picture of how they function on the data these models were trained to process.

Future work should extend this analysis to larger models, examine whether naturalistic-only induction heads emerge with scale, track how the QK--OV dissociation evolves during training using Pythia checkpoints, and connect our OV copying findings to the semantic induction mechanisms described by \citet{ren2024identifying}.
